% sec-01: the $200 billion realignment and the individual scientist.
% Resolved from science-golden-age chunk-01 and chunk-03.  Figures 1 and 2.
\section{The \$200 Billion Realignment and the Individual Scientist}

\subsection{What the Report Actually Says}

The July 2026 report to the President, \textit{Science: A New Golden Age},
states four overarching goals \cite{kratsios2026sciencegoldenage}: prioritise
the individual scientist over legacy institutions; change how research dollars
are allocated, distributed and assessed; set clear scientific goals and build
capacity to translate them; and prepare the enterprise for the AI revolution.

The first is stated without hedging: ``If we are serious about expanding what
our scientists can do, we must invest directly in American researchers and the
bold ideas that drive them. Too much of our research enterprise has come to
serve itself rather than the scientists within it.'' The Summary turns that into
an instruction with a named model attached: ``bet on people, not just
projects,'' by ``scaling long-horizon grants for the best and brightest modeled
on National Institutes of Health (NIH) Director's Pioneer Award.''

Chapter II attaches a number. Federal agencies distribute approximately
\textbf{\$200 billion} in annual research and development funding, and the
finding about that portfolio is the sentence this paper is written against:
there is ``no systematic framework for identifying where those dollars could
catalyze the greatest scientific returns, with deference instead given to the
same incumbents that consume the funding.''

The corrective has three parts: realign toward risk taking, strip away the
administrative burden that consumes close to half of an investigator's working
hours, and embed continuous evidence-based improvement. What the report does not
do is say who should receive the realigned money. This paper is one answer,
submitted ten times.

\begin{appfloat}[!tb]
\begin{appfig}[mermaid-type / flowchart]
% Four ranks.  Policy at y = 0, portfolio at y = -1.9, three mechanism families
% at y = -3.9 on a 4.9 pitch, ten recipients at y = -6.0 on a 2.35 pitch inside
% each family.  Families are 4.9 apart, so no leaf can touch a neighbouring
% family's leaf and no edge crosses a node.
\node[mmgoal,text width=52mm] (pol) at (5.5,0)
  {\textit{Science: A New Golden Age}, July 2026\\prioritize the individual
   scientist over legacy institutions};
\node[mmmid,text width=52mm] (b200) at (5.5,-1.9)
  {approximately \$200 billion annual federal R\&D\\``no systematic framework''};
\draw[mmedgeb] (pol) -- (b200);

\node[mmsoft,text width=30mm] (per) at (0.6,-3.9) {Person-based,\\long horizon};
\node[mmsoft,text width=30mm] (org) at (5.5,-3.9) {Organization-type,\\novel performer};
\node[mmsoft,text width=30mm] (par) at (10.4,-3.9) {Partnership\\and cost share};
\draw[mmedgeb] (b200) -- (per);
\draw[mmedgeb] (b200) -- (org);
\draw[mmedgeb] (b200) -- (par);

\node[mmin,text width=21mm] (a01) at (-0.6,-6.0) {01 NIH Pioneer};
\node[mmin,text width=21mm] (a07) at (1.8,-6.0) {07 HHMI};
\node[mmin,text width=21mm] (a02) at (2.2,-7.4) {02 ARPA-H};
\node[mmin,text width=21mm] (a03) at (4.6,-6.0) {03 X-Labs};
\node[mmin,text width=21mm] (a09) at (7.0,-6.0) {09 FRO};
\node[mmin,text width=21mm] (a05) at (4.6,-7.4) {05 SBIR};
\node[mmin,text width=21mm] (a04) at (9.2,-6.0) {04 Genesis};
\node[mmin,text width=21mm] (a06) at (11.6,-6.0) {06 FNIH};
\node[mmin,text width=21mm] (a08) at (9.2,-7.4) {08 NCI CTEP};
\node[mmgray,text width=21mm] (a10) at (11.6,-7.4) {10 UCSD Moores};
\foreach \s/\t in {per/a01, per/a07, org/a03, org/a09, par/a04, par/a06}{
  \draw[mmedge] (\s) -- (\t);}
\draw[mmedge] (org.south) to[out=-120,in=90,looseness=0.8] (a02.north);
\draw[mmedge] (org.south) to[out=-90,in=90,looseness=0.8] (a05.north);
\draw[mmedge] (par.south) to[out=-110,in=90,looseness=0.8] (a08.north);
\draw[mmedge] (par.south) to[out=-70,in=90,looseness=0.8] (a10.north);
\node[font=\tiny\itshape,text=protogray,anchor=north,text width=112mm,align=center]
  at (5.5,-8.2) {Recipient 10 is tinted because it is the only one that is not a
  funder: it is the partner institution, and what is asked of it is a
  feasibility review rather than money.};
\end{appfig}
\vspace{-0.7cm}
\figcaption{One sentence of federal policy, three mechanism families, and the
ten recipients each family reaches. The split is by mechanism, not by agency,
which is why one recipient is not a federal funder.}
\end{appfloat}

\subsection{Where the Enterprise Stalls}

The report's diagnosis is quantitative, and every number below is the report's
own or is drawn from the primary source its end notes name.

\begin{apptable}
\begin{tabularx}{\textwidth}{>{\raggedright\arraybackslash}p{5.0cm} >{\raggedright\arraybackslash}p{2.4cm} >{\raggedright\arraybackslash}X}
\headrow \thead{Finding} & \thead{Value} & \thead{Mechanism the report names against it}\\
\midrule
Research time lost to administration & near one half & Burden reduction; applications of a few pages \\
New federal requirements on grants, 1991 to January 2025 & at least 270 & Requirement rescission; metascience units \\
Negotiated indirect cost rates & 50 to 60 percent & Indirect-cost reform; private funders accept 10 to 15 \\
Grant lead time & up to 20 months & Fast grants decided in under one month \\
NIH principal investigator average age, 1980 to 2008 & 39 rising to 51 & Early-career independence; portable fellowships \\
Fall in new drugs per billion dollars since 1950 & roughly eightyfold & Portfolio construction across exploration and exploitation \cite{scannell2012eroom} \\
Publication-rate advantage of person-based funding & nearly double & HHMI-style support, roughly \$10M over seven years \cite{azoulay2011incentives} \\
\end{tabularx}
\end{apptable}
\tabcap{The report's diagnosis in seven numbers, each paired with the mechanism
it names as the corrective. The last row is the only one that reports an
outcome rather than a friction.}

An unaffiliated scientist meets a subset of these frictions and one the table
does not contain: the eligibility screen. Most federal mechanisms require a host
institution before a proposal can be reviewed at all, so the merit question is
never reached. Figure 2 states the resulting state machine and names, for each
stall, the mechanism the report offers against it.

\begin{appfloat}[!tb]
\begin{appfig}[plantuml-type / state diagram with guards]
% Four progress states on the upper rank at y = 0, pitch 3.4.  Three blocked
% states on the lower rank at y = -2.5, each directly below the state it
% interrupts, so the pairing is read from vertical alignment rather than from
% following an edge.  Return curves use bend left=16 with stated looseness.
\node[umlinit] (i) at (-1.6,0) {};
\node[umlstate,text width=24mm] (s1) at (0.2,0) {Drafted};
\node[umlstate,text width=24mm] (s2) at (3.6,0) {Submitted};
\node[umlstate,text width=24mm] (s3) at (7.0,0) {Reviewed};
\node[umlstateon,text width=24mm] (s4) at (10.4,0) {Funded};
\node[umlfinal] (f) at (12.2,0) {};
\draw[umlarrow] (i) -- (s1);
\draw[umlarrow] (s1) -- node[umlguard,pos=0.5] {[eligible host]} (s2);
\draw[umlarrow] (s2) -- node[umlguard,pos=0.5] {[passes screen]} (s3);
\draw[umlarrow] (s3) -- node[umlguard,pos=0.5] {[risk tolerated]} (s4);
\draw[umlarrow] (s4) -- (f);

\node[umlstategray,text width=27mm] (b1) at (0.2,-2.5) {Blocked: no degree-granting host};
\node[umlstategray,text width=27mm] (b2) at (3.6,-2.5) {Blocked: 270 added requirements};
\node[umlstategray,text width=27mm] (b3) at (7.0,-2.5) {Blocked: panel consensus favours the incumbent};
\draw[umlarrow] (s1.south) -- (b1.north);
\draw[umlarrow] (s2.south) -- (b2.north);
\draw[umlarrow] (s3.south) -- (b3.north);
\draw[umldash] (b1.east) to[out=20,in=-90,looseness=0.9]
  node[umlguard,pos=0.55] {X-Labs, FRO, SBIR} (s1.south east);
\draw[umldash] (b2.east) to[out=20,in=-90,looseness=0.9]
  node[umlguard,pos=0.55] {fast grants} (s2.south east);
\draw[umldash] (b3.east) to[out=20,in=-90,looseness=0.9]
  node[umlguard,pos=0.55] {Pioneer, HHMI} (s3.south east);
\end{appfig}
\vspace{-0.7cm}
\figcaption{The five states a proposal from an unaffiliated scientist passes
through, and the three at which the report says the enterprise currently stalls.
Each stall names the mechanism that removes it.}
\end{appfloat}

\subsection{Why an Independent Scientist Is the Test Case}

Between June 2025 and July 2026, one person produced fourteen deposited works:
two clinical protocols, an investigational new drug application, two funding
applications, five legislative drafts, and four guidance works
\cite{foundingdocs}. None of it carried indirect cost. Three quantitative trial
simulations were completed at an estimated \$36,330 each against an industry
benchmark that runs from \$120,000 for a single empirical virtual trial to more
than \$2 million for a ten-arm quantitative systems pharmacology trial
\cite{daraxsims}.

The claim is a ratio and nothing more. Not that the work is better than an
institution's, and not that institutions are unnecessary: only that the part of
the pipeline which used to require a program office no longer does, and that the
part which still requires an institution, patients and a theatre and a review
board, is exactly the part these ten applications are asking to reach.

\begin{apptable}
\begin{tabularx}{\textwidth}{>{\raggedright\arraybackslash}p{4.6cm} >{\raggedright\arraybackslash}X}
\headrow \thead{Transmittal-letter goal} & \thead{What the ten applications supply against it}\\
\midrule
Prioritise the individual scientist over legacy institutions & One named applicant, no host institution, fourteen deposited works before any award \\
Change how research dollars are allocated and assessed & Ten mechanism families addressed on the same science, so the mechanism is the only variable \\
Set clear goals and build translational capacity & A published protocol and IND, and a governance package designed to be site-portable \\
Prepare the enterprise for the AI revolution & An advisory system whose boundary is measured in milliseconds and released as an open schema \\
\end{tabularx}
\end{apptable}
\tabcap{The report's four overarching goals, and the specific thing the
application set supplies against each. The second row is what makes the set a
comparison rather than ten separate requests.}
